\subsection{\bt{Repair and its boundary}{修复及其边界}}
\label{sec:repair-eval}
\bi{\emph{The reference decides (Q3).} The last row of Table~\ref{tab:replay} repeats \system\ with a second G8 reference: instead of the benchmark's judged learning query, which was written knowing the status, version, and current-flag columns, it uses the SQL the agent wrote for the learning question, the reference a deployment would actually have. Over the same \RpPairN\ questions, correct answers fall from \RpPairJudge\ to \RpPairEx, \RpRefDiff\ points (95\% cluster-bootstrap interval \RpRefLo--\RpRefHi). Every lost answer becomes unavailable and none becomes wrong. Repairs of versioned corrections now succeed for \RpRevisionEx\ instead of \RpRevisionJudge\ of \RpRevisionN\ questions, and those of dimension history for \RpDimhistEx\ instead of \RpDimhistJudge\ of \RpDimhistN, because these changes touch the learning period and the agent's own query ignores the column that separates the versions; status-history repairs are unaffected (\RpStatusEx\ of \RpStatusN) because that change leaves the learning period intact. The successful status and versioning repairs reported in Section~\ref{sec:scenarios} therefore rely on references that anticipate those columns.}
{\emph{参照决定成败（Q3）。}表~\ref{tab:replay} 最后一行用另一种 G8 参照重复 \system：不用基准中经过判题的学习查询（它在编写时已知道状态、版本和当前标志列），而用智能体为学习题写下的 SQL，这才是部署中真正可得的参照。在同一批 \RpPairN\ 道题上，答对数从 \RpPairJudge\ 降到 \RpPairEx，下降 \RpRefDiff\ 个百分点（按库整群自助法 95\% 区间 \RpRefLo--\RpRefHi）。丢失的答案全部变为不可用，没有一个变为错误。版本化更正的修复现在只在 \RpRevisionN\ 道题中的 \RpRevisionEx\ 道上成功（原为 \RpRevisionJudge\ 道），维表拉链只在 \RpDimhistN\ 道中的 \RpDimhistEx\ 道上成功（原为 \RpDimhistJudge\ 道），因为这两种变化触及学习期，而智能体自己的查询不涉及区分版本的列；状态流水的修复不受影响（\RpStatusN\ 道中 \RpStatusEx\ 道），因为该变化没有触及学习期。因此，第~\ref{sec:scenarios} 节报告的状态流水与版本化更正修复的成功，依赖于预先知道这些列的参照。}

\bi{\emph{Ambiguous fixes.} The backup-copy change corrects some 2002 return amounts and appends a full backup copy of the returns table, marked by a source column, that keeps the old amounts. Filters on the primary source and on the backup source each restore one row per return without losing returns. Without the uniqueness rule, repair takes the first filter it finds, the backup. A reference that ignores the source column rejects it, but a reference that already filters on the primary source accepts it, and \system\ then publishes a revision that answers two of four new questions wrongly for both return metrics. With the uniqueness rule, repair finds both filters and withdraws the definitions under either reference. Rerunning the replay of \system\ with the rule changes none of the 3,000 paired outcomes, because the status, versioning, and dimension-history fixes are unique.}
{\emph{有歧义的修复。}备份副本变化更正了部分 2002 年的退货金额，并追加一整份退货表的备份副本（以来源列标记），副本保留旧金额。按主数据来源过滤和按备份来源过滤，都能使每笔退货恰好一行且不丢退货。没有唯一性规则时，修复取搜索到的第一个过滤，即备份。不涉及来源列的参照会拒绝它，但预先按主数据来源过滤的参照会接受它，于是 \system\ 发布的修订在两个退货指标的 4 道新题中都答错 2 道。有了唯一性规则，修复会找到两个过滤，在任一参照下都撤下定义。打开该规则后重新回放 \system，3,000 道配对题的结果没有任何变化，因为状态流水、版本化更正与维表拉链的修复都是唯一的。}

\bi{\emph{Availability during repair.} Rejecting requests that arrive while a repair runs makes \WrDefOff\ of \WrUses\ simultaneous uses (eight agents, 19 definitions, three runs at 200,000 rows) unavailable under definition-level and \WrCondOff\ under condition-level maintenance, where shared verdicts revoke many definitions at once; waiting for the in-flight repair removes all of them at the same total waiting time and database work.}
{\emph{修复期间的可用性。}若拒绝修复进行中到达的请求，同时到达的 \WrUses\ 次使用（8 个智能体、19 个定义、20 万行、3 轮）中，定义级维护有 \WrDefOff\ 次不可用，条件级维护有 \WrCondOff\ 次——共享结论使许多定义同时被撤销；让这些请求等待进行中的修复后，不可用全部消失，总等待时间与数据库工作量不变。}

